150c150
< \footnotetext{\dag~Supplementary Information available: detailed equations and analytical workflow. See DOI: 00.0000/00000000.}
---
> \footnotetext{\dag~Supplementary Information available: detailed equations, raw data and analytical workflow. See DOI: 00.0000/00000000.}
162c162
< Here we introduce a sustainable two-stage analytical strategy based on Gran and Schwartz linearization methods for significantly more reagent-efficient equivalence point determination. 
---
> Here we present \textsf{GranTED}, an open-source Python toolkit that implements an automated, two-stage Gran--Schwartz linearization strategy for significantly more reagent-efficient equivalence point determination. 
164,165c164,165
< In the first stage, a ordinary, full titration is performed past the equivalence point to establish a reliable reference equivalence point (EQP*, full curve).
< In the second stage, systematic backward trimming of the same dataset combined with Gran--Schwartz linearization identifies the earliest acceptable volume ($V_{\rm opt}$) at which the Gran--Schwartz equivalence point (EQP, shorten curve) converges within a defined accuracy threshold.
---
> In the first stage, a conventioal, full titration is performed past the equivalence point to establish a reliable reference equivalence point (EQP*).
> In the second stage, systematic backward trimming of the same dataset combined with Gran--Schwartz linearization identifies the earliest acceptable volume ($V_{\rm opt}$) at which the Gran--Schwartz equivalence point (EQP) converges within a defined accuracy threshold (0.010--0.050\,mL).
169c169
< Depending on the chosen accuracy threshold, e.g. 0.010--0.050\,mL, reagent savings of approximately 10--30\,\% per titration were achieved while maintaining excellent agreement with the reference EQP* on replicated determinations.
---
> Depending on the chosen accuracy threshold, reagent savings of approximately 10--30\% per titration were achieved while maintaining excellent agreement with the reference EQP*.
172c172
< Implemented in the open-source Python tool \textsf{GranTED}, this analytical strategy is freely available, providing an accessible, reproducible, and sustainable platform for both research, quality control and teaching laboratories.}}
---
> \textsf{GranTED} is freely available, providing an accessible, reproducible, and sustainable platform for both research, quality control and teaching laboratories.}}
182,183c182,183
< Potentiometric titrations are among the most widely used analytical techniques in both research, quality control and teaching laboratories for the accurate determination of acid--base equilibria, redox processes, complexation and precipitation reactions.
< Conventional procedures typically require titration well beyond the equivalence point (post-EQP, full-curve) to reliably identify the inflection point via derivative analysis.
---
> Potentiometric titrations are among the most widely used analytical techniques in both research, quality control and teaching laboratories for the accurate determination of acid--base equilibria, redox processes, and complexation reactions.
> Conventional procedures typically require titration well beyond the equivalence point (post-EQP) to reliably identify the inflection point via derivative analysis.
186c186
< Classical linearization methods, first introduced by Gunnar Gran in 1950--1952\,\cite{Gran1950,Gran1952} and later extended by Schwartz in 1987\,\cite{Schwartz1987}, offer a powerful alternative by transforming the sigmoidal titration curve into straight lines whose intersection to the volume axis directly yields the equivalence point.
---
> Classical linearization methods, first introduced by Gunnar Gran in 1950--1952 and later extended by Schwartz in 1987, offer a powerful alternative by transforming the sigmoidal titration curve into straight lines whose intersection directly yields the equivalence point.
190c190
< Here we introduce analytical strategy \textsf{GranTED} (Gran Titration Equivalence point Determination), available as an an open-source Python toolkit at \href{https://github.com/sgiani95/GranTED/}{\texttt{https://github.com/sgiani95/GranTED/}}, that fully automates Gran and Schwartz linearization while incorporating a novel two-stage optimisation approach.
---
> Here we introduce analytical strategy \textsf{GranTED} (Gran Titration Equivalence point Determination), available as an an open-source Python toolkit that fully automates Gran and Schwartz linearization while incorporating a novel two-stage optimisation approach.
192c192
< In the second stage, systematic backward trimming of the same dataset identifies the earliest acceptable volume (\(V_{\rm opt}\)) at which the Gran--Schwartz equivalence point (EQP) diverges from EQP* within a user-defined accuracy threshold.
---
> In the second stage, systematic backward trimming of the same dataset identifies the earliest acceptable volume (\(V_{\rm opt}\)) at which the Gran--Schwartz equivalence point (EQP) converges within a user-defined accuracy threshold.
203,229c203,210
< Tris(hydroxymethyl)aminomethane (\ce{TRIS},
< %
< \colorbox{orange}{$\geq$\,99.7\,\%}
< %
< ) and potassium hydrogen phthalate (\ce{KHP},
< %
< \colorbox{orange}{$\geq$\,99.7\,\%}
< %
< ) were obtained from Merck Sigma-Aldrich and used as received.
< Hydrochloric acid (0.1\,M Titripac certified standard) and sodium hydroxide (0.1\,M Titripac certified standard) solutions were purchased from Merck Millipore.
< Ultra-pure water (18.2 M$\Omega\,\cdot\,{\rm cm}$) was used for all solution preparations and dilutions.
< 
< Potentiometric measurements were carried out with a Mettler-Toledo T5 automatic titrator equipped with a combined glass pH electrode (
< %
< \colorbox{orange}{DGi111-SC}
< %
< ).
< Titrant delivery was performed with a
< %
< \colorbox{orange}{20\,mL}
< %
< burette (0.001\,mL resolution).
< All experiments were conducted at room temperature under magnetic stirring (
< %
< \colorbox{orange}{300\,rpm}
< %
< ).
---
> Tris(hydroxymethyl)aminomethane (TRIS, $\geq$99.8\,\%) and potassium hydrogen phthalate (KHP, $\geq$99.95\,\%) were obtained from Sigma-Aldrich and used as received.
> Hydrochloric acid (0.1\,M certified standard) and sodium hydroxide (0.1\,M certified standard) solutions were purchased from Merck.
> Ultra-pure water (18.2\,M$\Omega$\,cm) was used for all solution preparations and dilutions.
> 
> Potentiometric measurements were carried out with a Metrohm 905 Titrando automatic titrator equipped with a combined glass pH electrode (Metrohm Unitrode).
> The electrode was calibrated daily using commercial pH 4.00, 7.00, and 10.00 buffers.
> Titrant delivery was performed with a 20\,mL burette (0.001\,mL resolution).
> All experiments were conducted at $25.0 \pm 0.1$\,°C under magnetic stirring (300\,rpm).
234c215
< Two aqueous solution standard benchmark systems were investigated: (i) \ce{TRIS}/\ce{HCl} (weak base–strong acid) and (ii) \ce{KHP}/\ce{NaOH} (weak acid–strong base).
---
> Two standard benchmark systems were investigated: (i) \ce{TRIS}/\ce{HCl} (weak base–strong acid) and (ii) \ce{KHP}/\ce{NaOH} (weak acid–strong base).
236,238c217,219
< For all titrations, 5.00\,mL of 0.1\,M \ce{TRIS} or \ce{KHP} solution were diluted with 60\,mL deionized water and titrated with 0.1\,M \ce{HCl} or \ce{NaOH} respectively, using fixed 0.05\,mL increments.
< In both cases, the titration was continued well past the equivalence point %(typically to 120–150\,\% of the expected equivalence volume)
< to generate the complete dataset required for reference EQP* determination. Each system was analysed in sextuplicate.
---
> For \ce{TRIS}/\ce{HCl} titrations, 50.00\,mL of approximately 0.05\,M \ce{TRIS} solution was titrated with 0.1\,M \ce{HCl}.
> For \ce{KHP}/\ce{NaOH} titrations, 50.00\,mL of approximately 0.05\,M \ce{KHP} solution was titrated with 0.1\,M \ce{NaOH}.
> In both cases, the titration was continued well past the equivalence point (typically to 120–150\,\% of the expected equivalence volume) to generate the complete dataset required for reference EQP* determination. Each system was analysed in sextuplicate.
242c223
< All raw titration data (volume vs.\ potential) were processed using the developed open-source Python toolkit \textsf{GranTED} (version 2.0.21). The software automatically converts potential to pH, computes both Gran (\(g_1\)) and Schwartz (\(g_s\)) functions, identifies optimal linear regions via derivative-based segmentation and R$^2$ maximisation, and performs the backward trimming analysis at the core of the algorithm. User-defined parameters were different EQP deviation threshold 0.010\,mL, 0.020\,mL and 0.050\,mL. All calculations were performed on a standard laptop (Python 3.11 environment).
---
> All raw titration data (volume vs.\ potential) were processed using the developed open-source Python toolkit \textsf{GranTED} (version 1.0.1). The software automatically converts potential to pH, computes both Gran (\(g_1\)) and Schwartz (\(g_s\)) functions, identifies optimal linear regions via derivative-based segmentation and $R^2$ maximisation, and performs the backward trimming analysis at the core of the algorithm. User-defined parameters were different EQP deviation threshold 0.010--0.050\,mL. All calculations were performed on a standard laptop (Python 3.11 environment).
244c225
< The same datasets were also analysed with the conventional inflection-point method (first and second derivative) using the built-in Mettler-Toledo software for direct comparison.
---
> The same datasets were also analysed with the conventional inflection-point method (first and second derivative) using the built-in Mettler software for direct comparison.
248c229
< \subsection*{Method Development: Earliest Acceptable Volume and Gran--Schwartz Optimization}
---
> \subsection*{Method Development: Earliest Acceptable Point and Gran--Schwartz Optimization}
250c231
< The key innovation of the analytical strategy is a two-stage procedure that combines classical Gran and Schwartz linearization with systematic backward data trimming to identify the earliest volume at which a reliable equivalence point can be obtained. 
---
> The core innovation of the analytical strategy is a two-stage procedure that combines classical Gran and Schwartz linearization with systematic backward data trimming to identify the earliest volume at which a reliable equivalence point can be obtained. 
252c233
< In the first stage, a conventional full titration is performed well past the equivalence point.
---
> In the first stage, a conventional full titration is performed well past the equivalence point (post-EQP).
257c238
< Starting from the full titration and progressively removing data points from the post-EQP* region, the Gran--Schwartz linearization is reapplied at each step.
---
> Starting from the full titration and progressively removing data points from the post-EQP region, the Gran--Schwartz linearization is reapplied at each step.
259c240
< The difference $|\text{EQP} - \text{EQP*}|$ is evaluated against a defined accuracy threshold (e.g. 0.010--0.050\,mL). 
---
> The difference $|\text{EQP} - \text{EQP*}|$ is evaluated against a defined accuracy threshold (0.010--0.050\,mL). 
264,265c245,246
< Figure~\ref{fig:titrations} a) shows a representative conventional titration curve and Figure~\ref{fig:titrations} b) the corresponding Gran--Schwartz plot for the full dataset and Figure~\ref{fig:titrations} c) and d) the trimmed dataset at $V_{\rm opt}$ for the TRIS/HCl system.
< A convergence table (Table~\ref{tab:trimming}) illustrates how R$^2$, EQP, its standard error and the difference from reference EQP* evolve with the number of retained data points, clearly identifying the earliest acceptable volume.
---
> Figure~\ref{fig:titrations}a) shows a representative titration curve and Figure~\ref{fig:titrations}b) the corresponding Gran--Schwartz plot for the full dataset and Figure~\ref{fig:titrations}c) and d) the trimmed dataset at $V_{\rm opt}$ for the TRIS/HCl system.
> A convergence table (Table~\ref{tab:trimming}) illustrates how R$^2$, EP, and its standard error evolve with the number of retained data points, clearly identifying the earliest acceptable volume.
271,275c252,253
<   \caption{Overview of the green analytical strategy for equivalence point (EQP) determination in potentiometric titrations, demonstrated for TRIS (0.1 M) titrated with \ce{HCl} (0.1 M).
<   Left: a) Raw pH vs. titrant volume ($V$) curve from the initial titration (full run to post-EQP*). %, identifying the pre-EQP linear region (shaded) and its central truncation point (arrow).
<   b) Schwartz--Gran modified plot ($V \times 10^{-\mathrm{pH}}$ vs. $V$), showing the straight pre-EQP* line extrapolated to the EQP* (dashed line; intercept at 5.005\,mL).
<   Right: c) Titration performed up to the optimal earliest acceptable point volume (4.40\,mL in this case).
<   d) Corresponding Schwartz--Green modified plot confirming EQP extrapolation (intercept at 5.002\,mL) with reduced dispensed titrant volume.}
---
>   \caption{Overview of the green analytical strategy for equivalence point (EQP) determination in potentiometric titrations, demonstrated for TRIS (0.1 M) titrated with \ce{HCl} (0.1 M). Left: a) Raw pH vs. titrant volume ($V$) curve from the initial titration (full run to post-EQP). %, identifying the pre-EQP linear region (shaded) and its central truncation point (arrow).
>   b) Schwartz--Gran modified plot ($V \times 10^{-\mathrm{pH}}$ vs. $V$), showing the straight pre-EQP line extrapolated to the EQP (dashed line; intercept at 5.020\,mL). Right: c) Titration performed up to the optimal earliest acceptable point volume. d) Corresponding Schwartz--Green modified plot confirming EQP extrapolation (intercept at 5.002\,mL) with reduced dispensed volume.}
281,284c259
<   \caption{Gran--Schwartz results for the \ce{TRIS}/\ce{HCl} system in the vicinity of the optimal volume $V_{\rm opt}$ (4.40\,mL, in bold).
<   Full titration EQP* corresponds to 5.005\,mL.
<   The selected volume meets the acceptance criteria of deviation from reference EQP* smaller than 0.010\,mL) with excellent linearity and low uncertainty.
<   Except for R$^2$, all values in\,mL}
---
>   \caption{Gran--Schwartz results for the \ce{TRIS}/\ce{HCl} system in the vicinity of the optimal volume $V_{\rm opt}$ (4.40\,mL, in bold). Full titration EQP* corresponds to 5.007\,mL. The selected volume meets the acceptance criteria of deviation from reference EQP* smaller than 0.010\,mL) with excellent linearity and low uncertainty. Except for $R^2$, all values in\,mL}
288c263
<     Dispensed &     & Standard &       & Difference from \\
---
>     Dispensed &     & Standard &       & Deviation from \\
291,297c266,272
<     4.25 & 4.980 & 0.010 & 0.9987 & -0.025 \\
<     4.30 & 4.987 & 0.010 & 0.9988 & -0.018 \\
<     4.35 & 4.989 & 0.009 & 0.9989 & -0.016 \\
<     \textbf{4.40} & \textbf{5.002} & \textbf{0.005} & \textbf{0.9997} & \textbf{-0.003} \\
<     4.45 & 5.002 & 0.005 & 0.9997 & -0.003 \\
<     4.50 & 5.007 & 0.004 & 0.9998 & \phantom{-}0.002 \\
<     4.55 & 5.006 & 0.003 & 0.9998 & \phantom{-}0.001 \\
---
>     4.25 & 4.980 & 0.010 & 0.9987 & -0.027 \\
>     4.30 & 4.987 & 0.010 & 0.9988 & -0.020 \\
>     4.35 & 4.989 & 0.009 & 0.9989 & -0.018 \\
>     \textbf{4.40} & \textbf{5.002} & \textbf{0.005} & \textbf{0.9997} & \textbf{-0.005} \\
>     4.45 & 5.002 & 0.005 & 0.9997 & -0.005 \\
>     4.50 & 5.007 & 0.004 & 0.9998 & \phantom{-}0.000 \\
>     4.55 & 5.006 & 0.003 & 0.9998 & -0.001 \\
302c277
< This trimming-based optimisation constitutes the central methodological advance, transforming a traditional full titration into a reagent-efficient procedure without requiring additional experiments.
---
> This trimming-based optimisation constitutes the key methodological advance, transforming a traditional full titration into a reagent-efficient procedure without requiring additional experiments.
310c285
< For both benchmark systems, the reference equivalence point EQP* (determined from the complete post-EQP* dataset) served as the ground-truth value.
---
> For both benchmark systems, the reference equivalence point EQP* (determined from the complete post-EQP dataset) served as the ground-truth value.
312,315c287,289
< In the \ce{TRIS}/\ce{HCl} system, the Gran--Schwartz EQP at \(V_{\rm opt}\) agreed with EQP* within the chosen threshold accuracy of 0.010--0.050\,mL.
< The conventional inflection-point method yielded comparable agreement but only after the full titration volume was reached.
< The EQP in this case corresponds to 5.007 $\pm$ 0.001\,mL.
< Similar performance was observed for the \ce{KHP}/\ce{NaOH} system, where the inflection-point EQP corresponds to 5.032 $\pm$ 0.001\,mL.
---
> In the TRIS/HCl system, the Gran--Schwartz EQP at \(V_{\rm opt}\) agreed with EQP* within the accuracy of 0.010--0.050\,mL (indeed the chosen threshold). The conventional inflection-point method yielded comparable agreement but only after the full titration volume was reached.
> The EQP in this case corresponds to 5.015 $\pm$ 0.001\,mL.
> Similar performance was observed for the KHP/NaOH system, where the inflection-point EQP corresponds to 5.032 $\pm$ 0.001\,mL.
319,320c293
<   \caption{Comparison of equivalence point (EQP*) determination by Gran--Schwartz linear interpolation on the full titration (post-EQP*) and Gran--Schwartz linear extrapolation stopped at the optimal optimised volume corresponding to the indicated deviation from reference EQP*.
<   Results are mean $\pm$ SD ($n=6$ replicates). All values are in\,mL}
---
>   \caption{Comparison of equivalence point (EQP*) determination by Gran--Schwartz linear interpolation on the full titration (post-EQP) and Gran--Schwartz linear extrapolation stopped at the optimal optimised volume corresponding to the indicated deviation from reference EQP*. Results are mean $\pm$ SD ($n=6$ replicates). All values are in\,mL}
327c300
<     \multirow{2}{3em}{\ce{TRIS}/\ce{HCl}} & \phantom{ $\pm$}5.003 &  \phantom{ $\pm$}4.990 &  \phantom{ $\pm$}5.005 &  \phantom{ $\pm$}4.992 \\
---
>     \multirow{2}{3em}{\ce{TRIS}/\ce{HCl}} & \phantom{ $\pm$}5.013 &  \phantom{ $\pm$}4.990 &  \phantom{ $\pm$}5.005 &  \phantom{ $\pm$}4.992 \\
336,337c309,310
< Importantly, the standard error (Table~\ref{tab:trimming}) and standard deviation (Table~\ref{tab:validation}) associated with the Gran--Schwartz EQP remained low and congruent to the chosen difference threshold.% to that of the full method reference EQP*, demonstrating that the early-stop strategy does not compromise precision.
< The comparison of the equivalence points obtained by the three approaches (reference EQP*, Gran--Schwartz EQP at \(V_{\rm opt}\) and conventional inflection point) across both systems confirms excellent consistency.
---
> Importantly, the standard error (Table~\ref{tab:trimming}) and standard deviation (Table~\ref{tab:validation}) associated with the Gran--Schwartz EQP remained low and comparable to that of the full method reference EQP*, demonstrating that the early-stop strategy does not compromise precision.
> The comparison of the equivalence points obtained by the three approaches (reference EQP*, Gran--Schwartz at \(V_{\rm opt}\) and conventional inflection point) across both systems confirms excellent consistency.
344c317
< The optimised-volume strategy implemented in the analytical strategy directly translates into measurable reductions in chemical consumption.
---
> The optimised-volume strategy implemented of the propose analytical strategy directly translates into measurable reductions in chemical consumption.
346c319
< \(V_{\rm opt}\) identified by the backward trimming procedure depended strongly on the chosen accuracy threshold. For the \ce{TRIS}/\ce{HCl} system, \(V_{\rm opt}\) values of 4.4\,mL, 4.3\,mL and 3.6\,mL were obtained at thresholds of 0.010\,mL, 0.020\,mL and 0.050\,mL, respectively. For the \ce{KHP}/\ce{NaOH} system, the corresponding optimised volumes were 4.5\,mL, 4.35\,mL and 3.2\,mL.
---
> (\(V_{\rm opt}\)) identified by the backward trimming procedure depended strongly on the chosen accuracy threshold. For the \ce{TRIS}/\ce{HCl} system, \(V_{\rm opt}\) values of 4.4\,mL, 4.3\,mL and 3.6\,mL were obtained at thresholds of 0.010\,mL, 0.020\,mL and 0.050\,mL, respectively. For the \ce{KHP}/\ce{NaOH} system, the corresponding optimised volumes were 4.5\,mL, 4.35\,mL and 3.2\,mL.
348c321
< These substantially reduced volumes translate directly into significant reagent savings. For both benchmark systems, reagent consumption decreased from approximately 10\,\% (at the strictest 0.010\,mL accuracy threshold) to approximately 30\,\% (at the 0.050\,mL threshold) relative to a conventional full titration continued well past the equivalence point.
---
> These substantially reduced volumes translate directly into significant reagent savings. For both benchmark systems, reagent consumption decreased from approximately 10\% (at the strictest 0.010\,mL accuracy threshold) to 30\% (at the 0.050\,mL threshold) relative to a conventional full titration continued well past the equivalence point.
350c323
< These savings were consistently achieved across both systems, as illustrated in Fig.~\ref{fig:volume_savings} as reagent-saving histogram.
---
> These savings were consistently achieved across both systems, as illustrated in Fig.~\ref{fig:volume_savings} (reagent-saving histogram).
355,356c328
<   \caption{Reagent savings achieved by stopping at the optimal volume using different acceptance criteria (0.010, 0.020 and 0.050\,mL difference from the full-dataset EQP*).
<   Results are shown for both \ce{TRIS}/\ce{HCl} and \ce{KHP}/\ce{NaOH} systems ($n=6$ replicates each) and compared against the equivalence point} % and full titration endpoint}
---
>   \caption{Reagent savings achieved by stopping at the optimal volume using different acceptance criteria (0.010, 0.020 and 0.050\,mL difference from the full-dataset EQP). Results are shown for both \ce{TRIS}/\ce{HCl} and \ce{KHP}/\ce{NaOH} systems ($n=6$ replicates each) and compared against the equivalence point and full titration endpoint}
361c333
< This approach avoids the unnecessary addition of excess titrant that is typical in conventional potentiometric procedures, which must extend well beyond the equivalence point to reliably identify the inflection point.
---
> This approach avoids the unnecessary addition of excess titrant that is typical in conventional potentiometric procedures, which must extend well beyond the equivalence point to reliably identify the inflection.
364,365c336,337
< Reducing titrant consumption lowers both the risk associated to reagents used and the volume of waste generated per analysis.
< Furthermore, the open-source nature of \textsf{GranTED} enables widespread adoption in academic and industrial laboratories, promoting more sustainable analytical practices without requiring commercial software.
---
> Reducing titrant consumption lowers both the amount of hazard of reagents used and the volume of waste generated per analysis.
> Furthermore, the open-source nature of \textsf{GranTED} enables widespread adoption in academic and industrial laboratories, promoting more sustainable analytical practices without requiring expensive commercial software.
373c345
< Applied to the benchmark systems \ce{TRIS}/\ce{HCl} and \ce{KHP}/\ce{NaOH}, this approach consistently delivered equivalence points in excellent agreement with both the reference EQP* and the conventional inflection-point method, while reducing titrant consumption by about 10--30\,\% per titration (depending on the set accuracy). The strategy requires no additional experiments or hardware and eliminates the subjectivity associated with manual linear-region selection.
---
> Applied to the benchmark systems \ce{TRIS}/\ce{HCl} and \ce{KHP}/\ce{NaOH}, this approach consistently delivered equivalence points in excellent agreement with both the reference EQP* and the conventional inflection-point method, while reducing titrant consumption by 10--30\% per titration. The strategy requires no additional experiments or hardware and eliminates the subjectivity associated with manual linear-region selection.
375c347
< \textsf{GranTED} therefore provides a practical, accurate, precise, and significantly greener alternative to traditional potentiometric titration workflows. By enabling earlier termination of titrations without loss of analytical quality, it directly supports Green Chemistry Principle 1 (prevention of waste) and Principle 5 (design of safer solvents and auxiliaries). As a freely available, MIT-licensed tool with comprehensive diagnostic and visualisation capabilities, \textsf{GranTED} is expected to find wide application in research, quality control, and teaching laboratories seeking more sustainable analytical practices.
---
> \textsf{GranTED} therefore provides a practical, accurate, and significantly greener alternative to traditional potentiometric titration workflows. By enabling earlier termination of titrations without loss of analytical quality, it directly supports Green Chemistry Principle 1 (prevention of waste) and Principle 5 (design of safer solvents and auxiliaries). As a freely available, MIT-licensed tool with comprehensive diagnostic and visualisation capabilities, GranTED is expected to find wide application in research, quality control, and teaching laboratories seeking more sustainable analytical practices.
393c365
< The datasets supporting this article are available at \href{https://github.com/sgiani95/GranTED/}{\texttt{https://github.com/sgiani95/GranTED/}} or from the corresponding authors upon reasonable request.
---
> The datasets supporting this article are available in the Electronic Supporting Information (ESI) or from the corresponding authors upon reasonable request.
